\chapter{Возвраты, неравновесный порядок и спины Изинга}

\lectureoverview{Газ, выпущенный из одной половины комнаты, быстро устремляется к равновесию, однако обратимая механика говорит, что при достаточно долгой эволюции должны происходить и возвраты в это особое состояние. Мы оценим время такого возврата по объёму фазового пространства, отделим это явление от необъяснённого низкоэнтропийного начала Вселенной и используем возникающую отсюда загадку, чтобы обсудить мозги Больцмана, согласованную историю и жизнь, поддерживаемую потоком энергии. Затем мы перейдём к магнитному порядку. Независимые спины с двумя состояниями дают точно решаемый подготовительный пример: комбинаторика определяет статистическую сумму, а намагниченность выражается гиперболическим тангенсом. Взаимодействия между соседями, наконец, восстанавливают симметрию вверх--вниз и приводят к одномерной модели Изинга.}

\section{Воздух в комнате, который в конце концов возвращается}

Начнём с того, что весь воздух находится в левой половине комнаты. Такое состояние можно подготовить с помощью перегородки: откачать воздух из правой половины, поместить его слева, а затем убрать перегородку. Газ быстро распространяется по комнате и приближается к тепловому равновесию. Энтропия возрастает. Ничего удивительного не произошло.

Теперь оставим комнату герметично закрытой и подождём. Миллион орлов подряд — чрезвычайно маловероятный результат серии подбрасываний монеты, но в эксперименте, повторяемом неограниченно долго, рано или поздно возникают и чрезвычайно маловероятные результаты. Та же логика применима и к газу. В какой-то очень далёкий момент все его молекулы снова окажутся в левой половине или в какой-нибудь столь же особой области. Такое возвращение — это возврат Пуанкаре.

Точная теорема сформулирована несколько аккуратнее, чем эта наглядная картина. Для области фазового пространства конечной меры при динамике, сохраняющей объём фазового пространства, почти всякая начальная точка возвращается сколь угодно близко к своему исходному положению. Макроскопическое утверждение вроде \emph{все молекулы находятся в левой половине} проще: оно требует лишь, чтобы траектория снова вошла в конечную огрублённую область, а не чтобы каждая молекулярная координата вернулась в точности. Практический вопрос состоит в том, сколько времени это занимает.

\subsection{Фазовая точка высокой размерности}

Для \(N\) классических частиц фазовое пространство имеет \(6N\) координат,
\begin{equation}
 \Gamma=(\bm x_1,\ldots,\bm x_N;\bm p_1,\ldots,\bm p_N).
 \label{eq:sm8-phase-point}
\end{equation}
Положения ограничены комнатой. При фиксированной полной энергии импульсы тоже ограничены, потому что одна частица не может нести сколь угодно большой импульс, не неся сколь угодно большой энергии. Поэтому допустимая энергетическая оболочка имеет конечную меру при тех предположениях, которые используются в рассуждении о возврате.

Фазовая точка движется детерминированно, но в многочастичном газе её движение хаотично. Близкие траектории быстро расходятся, как это происходит в последовательности столкновений бильярдных шаров. После огрубления длинная траектория кажется блуждающей по всей доступной области.

\begin{figure}[t]
\centering
\includegraphics[width=0.90\textwidth]{lecture_08_figure_02.png}
\caption{Схематическая область фазового пространства и сложная траектория, блуждающая по ней. Рисунок намеренно сделан низкоразмерным; реальный газ живёт в \(6N\)-мерном пространстве.}
\lectureframe{00:06:41}
\end{figure}

Низкоразмерный набросок сразу же подталкивает нас к ошибке. Если бы особое состояние занимало половину изображённого прямоугольника, можно было бы сказать, что газ должен проводить там половину времени. Но условие \emph{каждая молекула лежит в левой половине} накладывает ограничение на каждую координату положения, а не только на одну координату фазовой точки.

Возьмём всего две частицы, движущиеся в одном пространственном измерении, с координатами \(x_1\) и \(x_2\). Их конфигурационное пространство — квадрат. Условие \(x_1<L/2\) отбирает половину квадрата; условия \(x_1<L/2\) и \(x_2<L/2\) вместе отбирают одну четверть:
\begin{equation}
 P_2=\frac{1}{2}\frac{1}{2}=\frac{1}{4}.
 \label{eq:sm8-two-particle-quarter}
\end{equation}

\begin{figure}[t]
\centering
\includegraphics[width=0.90\textwidth]{lecture_08_figure_07.png}
\caption{Для двух одномерных частиц меньший квадрат представляет конфигурации, в которых обе координаты лежат в выбранной половине комнаты.}
\lectureframe{00:07:08}
\end{figure}

Для трёх частиц получаем \(1/2^3\), а для \(N\) частиц получаем
\begin{equation}
 P_{\mathrm{half}}(N)
 =\left(\frac{1}{2}\right)^N.
 \label{eq:sm8-half-probability}
\end{equation}
Весь смысл именно в показателе степени. Множитель одна вторая вполне обычен; одна вторая в степени макроскопического числа частиц выходит за пределы обычной интуиции.

\begin{classroomqa}
\audiencequestion{Разве ответ не должен быть \(1/3^N\), ведь у молекул есть и другие пространственные направления?}
\lecturerresponse{Это зависит от того, какое условие налагается. Если единственный вопрос состоит в том, лежит ли каждая частица слева от разделяющей стенки, то для каждой частицы это одно двоичное условие, и ответ равен \(2^{-N}\). Наложение ограничения ещё на одну независимую долю её доступного пространства положений дало бы ещё один множитель.
\lecturetimestamp{00:07:47--00:08:05}}
\end{classroomqa}

\section{Фазовый объём, энтропия и масштаб возврата}

Ту же оценку можно записать и без обращения к половинам. Пусть комната имеет обычный пространственный объём \(V\), а \(v\) — меньшая область, в которой мы требуем, чтобы находилась каждая молекула. Пренебрегая импульсными множителями, общими для двух макросостояний, получаем для доступных объёмов конфигурационного пространства
\begin{equation}
 \Omega_{\mathrm{room}}\propto V^N,
 \qquad
 \Omega_{\mathrm{small}}\propto v^N.
 \label{eq:sm8-config-volumes}
\end{equation}
Поэтому доля долгого времени, проводимая в малом макросостоянии, равна
\begin{equation}
 P(v)
 =\frac{\Omega_{\mathrm{small}}}{\Omega_{\mathrm{room}}}
 =\left(\frac{v}{V}\right)^N.
 \label{eq:sm8-volume-ratio}
\end{equation}

\begin{figure}[t]
\centering
\includegraphics[width=0.90\textwidth]{lecture_08_figure_08.png}
\caption{Общая схема фазовых объёмов на доске: малая область \(v^N\) внутри доступной области \(V^N\), с отношением \(v^N/V^N\).}
\lectureframe{00:10:44}
\end{figure}

В единицах, где \(k_B=1\), энтропия есть логарифм доступного фазового объёма. Если \(S_A=\log\Omega_A\) — энтропия ограниченного макросостояния, а \(S_{\rm eq}=\log\Omega_{\rm eq}\) — энтропия равновесия, то
\begin{equation}
 \begin{aligned}
 P(A)
 &=\frac{\Omega_A}{\Omega_{\rm eq}}\\
 &=\exp\!\left(S_A-S_{\rm eq}\right)\\
 &=e^{-\Delta S},\\
 \Delta S&=S_{\rm eq}-S_A.
 \end{aligned}
 \label{eq:sm8-entropy-probability}
\end{equation}
Для случая полкомнаты \(\Delta S=N\log 2\), и
\eqref{eq:sm8-entropy-probability} в точности сводится к \(2^{-N}\).
Сокращённая запись \(P\sim e^{-S_{\rm eq}}\) уместна только тогда, когда энтропия особой области рассматривается как фиксированный эталон; физически осмысленным показателем степени является дефицит энтропии \(\Delta S\).

\begin{classroomqa}
\audiencequestion{Это объём конфигурационного пространства, а не всего фазового пространства?}
\lecturerresponse{Для этого сравнения — да. Мы накладываем условие на положения, оставляя распределение по импульсам при той же температуре, поэтому общий множитель из пространства импульсов сокращается. Полностью общее утверждение о возврате использует фазовое пространство; эту конкретную вероятность можно вычислить из отношения объёмов конфигурационного пространства.
\lecturetimestamp{00:11:21--00:11:40}}
\end{classroomqa}

Предположим, что система теряет память о своём детальном положении за микроскопическое время корреляции или перемешивания \(\tau_{\rm mix}\). Тогда попадание в редкую область с вероятностью \(P(A)\) происходит на масштабе
\begin{equation}
 \tau_{\rm rec}
 \sim \frac{\tau_{\rm mix}}{P(A)}
 \sim \tau_{\rm mix}e^{\Delta S}.
 \label{eq:sm8-recurrence-time-general}
\end{equation}
Для условия полкомнаты
\begin{equation}
 \tau_{\rm rec}\sim\tau_{\rm mix}2^N.
 \label{eq:sm8-recurrence-time-half}
\end{equation}
В лекции используется намеренно грубая оценка для комнаты \(N\sim10^{30}\), дающая легко запоминающийся масштаб
\(\tau_{\rm rec}/\tau_{\rm mix}\sim2^{10^{30}}\). При обычной атмосферной плотности в комнате скорее \(10^{27}\) молекул, но замена \(10^{30}\) на \(10^{27}\) не меняет вывода: время возврата остаётся немыслимо большим по сравнению с возрастом Вселенной.

\begin{classroomqa}
\audiencequestion{В каких единицах измеряется время возврата? Может ли переход от секунд к часам или к чрезвычайно короткой единице времени изменить вывод?}
\lecturerresponse{Размерный префактор здесь — \(\tau_{\rm mix}\); гигантское число — это безразмерное отношение \(\tau_{\rm rec}/\tau_{\rm mix}\).
Переход от секунд к часам делит результат на обычный коэффициент перевода, который ничтожен рядом с \(2^{10^{30}}\). Длительность одного посещения угла действительно зависит от скоростей молекул. Оценка касается того, насколько редко начинается такое посещение. Существует и квантовая теорема о возврате, но здесь она не нужна.
\lecturetimestamp{00:17:24--00:19:18}}
\end{classroomqa}

\begin{classroomqa}
\audiencequestion{Какова цель вычисления времени, которого никакой эксперимент заведомо не сможет дождаться?}
\lecturerresponse{Смысл здесь концептуальный. Это показывает, в каком смысле замкнутая механическая система остаётся обратимой. Почти всё время она проводит вблизи равновесия, но на достаточно больших временных масштабах испытывает и флуктуации в сторону меньшей энтропии. Поэтому второй закон — это утверждение об особых начальных данных, которое верно с подавляющей вероятностью, а не микроскопическое уравнение, запрещающее обратное движение.
\lecturetimestamp{00:15:12--00:16:35}}
\end{classroomqa}

\section{Обратимость и особое начало}

Начнём сознательно с крошечной области, соответствующей половине комнаты. Почти всякая соседняя траектория немедленно покидает её, потому что снаружи находится несравнимо больший фазовый объём, чем внутри. Газ быстро распространяется, а затем проходит огромное время, прежде чем он вернётся. Эта асимметрия между быстрой релаксацией и долгим ожиданием не противоречит обратимой механике.

Обращённый по времени процесс быстрого выхода — не весь период ожидания. Это краткий, фантастически хорошо скоординированный эпизод, в котором молекулярные скорости, выглядящие равновесными, случайно выстраиваются так, что газ \emph{влетает} в одну половину комнаты. Если в конце обычного выхода обратить каждый импульс, газ повторит этот выход в обратном направлении. Долгое время ожидания — это время, необходимое для того, чтобы такая особая входящая конфигурация скоростей возникла спонтанно.

\begin{classroomqa}
\audiencequestion{Означает ли симметрия по времени, что медленный возврат в угол занимает столько же времени, сколько первоначальный выход из него?}
\lecturerresponse{Нет. Сравните выход с его точным обращением во времени. Оба процесса быстры. Огромное время занимает ожидание в равновесии той точной входящей флуктуации, которая запускает обратное стремительное движение. На чрезвычайно длинной временной записи отклонения от равновесия происходят в обоих направлениях времени с симметричными локальными историями.
\lecturetimestamp{00:20:46--00:21:58}}
\end{classroomqa}

Следовательно, аргумент о возврате объясняет лишь часть стрелы времени. Он объясняет, почему энтропия обычно возрастает \emph{при условии}, что система началась в области низкой энтропии. Он не объясняет, почему начальное условие было особым. В космологии эта недостающая посылка становится центральным вопросом: почему наблюдаемая Вселенная началась в такой исключительно малой области своего гравитационного фазового пространства?

\begin{classroomqa}
\audiencequestion{Применим ли возврат к космологии? После того как Вселенная расширится, может ли в ней при достаточно долгом ожидании произойти возврат?}
\lecturerresponse{Этот вопрос неизбежен, но ответ на него зависит от глобальной космологии. По-настоящему замкнутая система с конечным пространством равновесных состояний естественно допускает возврат. Расширяющееся пространство-время, гравитация, горизонты, а также бесконечное или нестационарное пространство состояний могут изменить предпосылки. Однако важный вывод здесь более узок: сама по себе картина равновесного возврата не является удовлетворительным статистическим объяснением той структурированной Вселенной, которую мы видим.
\lecturetimestamp{00:19:59--00:20:45}}
\end{classroomqa}

\begin{classroomqa}
\audiencequestion{Зависит ли возврат от того, что вначале газ подготовлен в малой области? Если вначале в комнате равновесие, соберётся ли воздух всё равно в конце концов в одном углу?}
\lecturerresponse{Да. Возврат относится к долговременной динамике замкнутой системы, поэтому и начальное состояние, выглядящее равновесным, тоже испытывает редкие низкоэнтропийные флуктуации. Подготовка состояния «полкомнаты» важна по другой причине: она позволяет нам рассматривать ситуацию при условии известного особого начала и спрашивать, почему сразу после него энтропия возрастает.
\lecturetimestamp{00:22:04--00:22:20}}
\end{classroomqa}

\section{Почему случайная флуктуация предсказывает слишком мало}

Можно попытаться превратить идею возврата в космологию: возможно, Вселенная провела почти всю вечность в равновесии, а наш нынешний мир — одна редкая флуктуация. Это предположение не выдерживает строгой статистической проверки. Рассмотрим ситуацию при условии, что наблюдатель существует, и спросим, что ещё такой наблюдатель должен ожидать увидеть.

Флуктуация, достаточно большая, чтобы создать одну планету, чрезвычайно редка. Флуктуация, создающая две планеты, требует большего дефицита энтропии и встречается ещё реже. Флуктуация, создающая звёзды, галактики, космический микроволновый фон и взаимно согласованные записи о долгом прошлом, встречается реже ещё на колоссальный дополнительный экспоненциальный множитель. Следовательно, теория равновесных флуктуаций предсказывает наименьшую флуктуацию, достаточную для наблюдения, а не тот мир с богатой внутренней согласованностью, который действительно наблюдается. В своей наиболее крайней форме минимальный наблюдатель называется мозгом Больцмана.

Логика здесь основана на условной вероятности. Пусть \(A\) — событие, состоящее в возникновении одной флуктуации масштаба наблюдателя, а \(B\) — событие существования второй независимой структуры. Если вероятность каждого мала и равна \(p\), то
\begin{equation}
 P(B\mid A)=p,
 \qquad
 P(A\cap B)=p^2.
 \label{eq:sm8-conditional-fluctuations}
\end{equation}
Тот факт, что первая флуктуация произошла, не делает вторую вероятной. Аналогия с лотереей на этом уровне точна: у предыдущего победителя нет повышенного шанса при следующем независимом розыгрыше. Поэтому наблюдатель может рассматривать ситуацию при условии собственного существования, но всё равно должен изумиться, обнаружив независимо структурированное окружение.

\begin{classroomqa}
\audiencequestion{Если одна планета или один больцмановский наблюдатель уже существует, действительно ли вероятность появления ещё одного меньше?}
\lecturerresponse{Условная вероятность второй независимой редкой флуктуации равна той же малой вероятности одного события, а не её квадрату. Совместная вероятность обеих — квадрат этой вероятности. Именно поэтому такой наблюдатель всё равно должен ожидать теплового беспорядка вне себя. Обнаружение супруга, другой планеты и упорядоченного неба служит свидетельством против предпосылки, что наблюдатель был собран как изолированная равновесная флуктуация.
\lecturetimestamp{00:26:12--00:30:27}}
\end{classroomqa}

\subsection{Согласованные записи о прошлом}

Тот же аргумент относится и к истории. Два учебника, которые согласуются друг с другом, астрономические записи, соответствующие одной истории расширения, геологические слои, воспоминания и лабораторные журналы — все они взаимно ограничивают друг друга. Один случайный учебник мог бы содержать любую историю. Целая сеть записей, согласующихся относительно одного и того же прошлого, занимает гораздо меньшую область фазового пространства.

В качестве примера на занятии берётся знакомая история о том, что Джордж Вашингтон срубил вишнёвое дерево. Истинна эта история или нет, несущественно. Существенно то, что в мире, возникшем флуктуацией в настоящее мгновение, предложение в одной книге не даёт никакой динамической причины тому, чтобы другая запись с ним согласовалась. Согласованные корреляции — это свидетельство общего прошлого, а не просто одной мгновенной случайности.

\begin{classroomqa}
\audiencequestion{Что здесь означает \emph{историческая согласованность}?}
\lecturerresponse{Это означает, что многие нынешние записи согласуются друг с другом как следствия общего прошлого. Флуктуация могла бы сфабриковать записи и воспоминания, но создание взаимно согласованной сети требует гораздо большей флуктуации, чем создание одного минимального наблюдателя. Поэтому теория возврата предсказывает, что большинству наблюдателеподобных флуктуаций недостаёт той согласованности, которую мы в действительности обнаруживаем.
\lecturetimestamp{00:30:30--00:31:12}}
\end{classroomqa}

\begin{classroomqa}
\audiencequestion{Может ли Вселенная всё же быть замкнутой системой, у которой времена установления равновесия и возврата просто чрезвычайно велики?}
\lecturerresponse{Такая возможность не снимает статистической проблемы. За сколь угодно долгую историю существовало бы много наблюдателеподобных флуктуаций, и подавляющее большинство из них были бы минимальными, без согласованной окружающей истории. Наблюдатель, выбранный из этого ансамбля, не должен был бы ожидать именно наших наблюдений. Поэтому космологи всерьёз рассматривают проблему мозга Больцмана как проверку космологических мер и моделей поздней эволюции, а не как утверждение, будто возврат уже наблюдался.
\lecturetimestamp{00:31:13--00:32:42}}
\end{classroomqa}

\section{Жизнь как структура, поддерживаемая потоком}

Жизнь, по-видимому, ставит противоположную трудность. Организмы поддерживают сложную организацию в течение долгого времени; они не похожи на краткие случайные флуктуации. Разрешение состоит в том, что второй закон относится к полной замкнутой системе. Подсистема может понижать свою энтропию или поддерживать её при условии, что в её окружении происходит ещё больший рост энтропии:
\begin{equation}
 \dot S_{\rm total}
 =\dot S_{\rm subsystem}+\dot S_{\rm environment}\geq 0.
 \label{eq:sm8-open-system-second-law}
\end{equation}

Земля не находится в тепловом равновесии. На человеческих временных масштабах она приблизительно стационарна, но стационарность не означает статичности. Через неё проходит поток энергии. Солнечный свет поступает, деградирует в ходе многих процессов и переизлучается при более низкой эффективной температуре. Такой сквозной поток может поддерживать локальную организацию, пока полная энтропия возрастает.

Наглядный механический образ даёт вода, текущая по трубе. Поток может содержать устойчивые завихрения и вихри. Это организованные структуры, но поддерживаются они течением. Если герметично закрыть оба конца и остановить поток, вязкость сотрёт завихрения по мере того, как вода будет приближаться к скучному равновесию. Вытекающая вода теплее входящей, так что более крупная система оплатила энтропийный счёт.

\begin{classroomqa}
\audiencequestion{Уменьшает ли жизнь энтропию и почему её организация может сохраняться, а не появляться и исчезать как обычная флуктуация?}
\lecturerresponse{Жизнь может локально поддерживать или понижать энтропию, но только экспортируя ещё больше энтропии. Её устойчивость обусловлена стационарным неравновесным потоком, прежде всего потоком солнечной энергии через Землю. Если бы Земля была изолирована от входящего и исходящего излучения, она в конце концов приблизилась бы к тепловому равновесию, и упорядоченная деятельность прекратилась бы. Жизнь ближе к вихрю в потоке, чем к равновесной флуктуации.
\lecturetimestamp{00:32:43--00:35:42}}
\end{classroomqa}

На этом мы оставим \emph{интересные вещи} и вернёмся, пользуясь нарочито сухой формулировкой лекции, к чему-то скучному: магнитам.

\section{Охлаждение магнита: домены и выбранное направление}

В статистической механике математическую модель часто называют магнитом, даже если это не буквальный кусок железа. Полезная общая черта состоит в наличии совокупности локальных степеней свободы, способных упорядочиваться. В реальном магнитном материале эти степени свободы могут возникать из атомных магнитных моментов, спинов электронов, орбитальных токов или эффективных моментов структурных единиц кристалла.

При высокой температуре тепловой беспорядок заставляет локальные моменты быть направленными во многие стороны, так что суммарная намагниченность мала или вовсе отсутствует. Если энергия соседних моментов ниже, когда они сонаправлены, то при охлаждении могут образовываться небольшие сонаправленные домены. При дальнейшем охлаждении энергетическое преимущество становится важнее, и коррелированные домены растут. В системе, допускающей ферромагнитный переход, длина корреляции становится макроскопической при конечной критической температуре. Тогда упорядоченная фаза выбирает некоторое направление, хотя энергия при нулевом поле заранее не предпочитала никакого определенного направления.

\begin{classroomqa}
\audiencequestion{Должны ли микроскопические магниты предпочитать антипараллельное выравнивание, поскольку северный полюс притягивает южный?}
\lecturerresponse{Эффективное взаимодействие зависит от материала.
Обычная дипольная интуиция не исчерпывает всей микроскопической задачи; важны обмен,
кристаллическая структура и конкурирующие взаимодействия. Ферромагнитные
материалы — это такие материалы, в которых эффективная низкоэнергетическая тенденция благоприятствует
параллельному порядку, тогда как в других материалах может предпочитаться антипараллельное выравнивание.
\lecturetimestamp{00:39:56--00:40:30}}
\end{classroomqa}

\begin{classroomqa}
\audiencequestion{Нужно ли, чтобы при охлаждении материала присутствовало внешнее магнитное поле, чтобы моменты выстроились?}
\lecturerresponse{Нет. Взаимодействие может породить порядок и без поля.
В точно симметричной идеализации разные направления энергетически
эквивалентны, поэтому реализующееся направление выбирается начальными флуктуациями,
границами или сколь угодно малым паразитным полем. Выбор одного из этих
эквивалентных направлений есть спонтанное нарушение симметрии; возникновение
макроскопического порядка — это магнитный фазовый переход.
\lecturetimestamp{00:40:31--00:42:33}}
\end{classroomqa}

Реальные материалы сложны, поэтому теперь мы пойдем на грубое упрощение. Каждый локальный
момент может быть направлен только вверх или вниз:
\begin{equation}
 \sigma_i=+1\quad\hbox{or}\quad\sigma_i=-1.
 \label{eq:sm8-ising-spin}
\end{equation}
В этой модели нет ни востока, ни запада. Два допустимых направления — это
определение модели, так же как орел и решка суть два исхода броска монеты.
Ее ценность не в микроскопическом реализме, а в математической ясности и способности
описывать явления упорядочения, общие для многих систем.

\begin{classroomqa}
\audiencequestion{Что физически представляют собой эти объекты с двумя состояниями? Следует ли
считать их частицами?}
\lecturerresponse{Можно представить себе атомы в узлах кристалла, каждый с
эффективным магнитным моментом, обусловленным спином или микроскопическим током. Для
модели, однако, важна только двузначная переменная. Ограничение
вверх и вниз вводится для того, чтобы статистическая механика оставалась простой.
\lecturetimestamp{00:44:22--00:45:01}}
\end{classroomqa}

\section{Независимые спины во внешнем поле}

Первая модель содержит \(N\) спинов, которые не взаимодействуют друг с другом.
Единственный энергетический масштаб задается внешним полем \(H\), а \(\mu\) характеризует
силу связи с ним одного спина. Сохраним то соглашение о знаках, которое использовалось
в вычислении:
\begin{equation}
 \epsilon_i=\mu H\sigma_i.
 \label{eq:sm8-single-spin-energy}
\end{equation}
При \(H>0\) спин вверх \(\sigma_i=+1\) имеет энергию \(+\mu H\), тогда как спин
вниз \(\sigma_i=-1\) имеет энергию \(-\mu H\). Следовательно, при этом соглашении
положительное поле благоприятствует спинам, направленным вниз. Это противоположно другому распространенному
соглашению в обозначениях, но физика не зависит от названий, если знаки используются
последовательно.

Пусть \(n\) — число спинов вверх, а \(m\) — число спинов вниз. Тогда
\begin{equation}
 n+m=N,
 \qquad
 E=(n-m)\mu H.
 \label{eq:sm8-counting-energy}
\end{equation}

\begin{figure}[t]
\centering
\includegraphics[width=0.90\textwidth]{lecture_08_figure_03.png}
\caption{Схема независимых спинов: \(n\) спинов вверх, \(m\) спинов вниз,
\(n+m=N\) и \(E=(n-m)\mu H\).}
\lectureframe{00:49:50}
\end{figure}

\begin{classroomqa}
\audiencequestion{Каковы размерности \(H\), \(\mu\) и \(E\)? Является ли
\(H\) числом или магнитным полем?}
\lecturerresponse{\(H\) обозначает величину внешнего приложенного поля, а
\(\mu\) — магнитный момент. Их произведение \(\mu H\) имеет размерность
энергии. В абстрактной модели их часто объединяют в один параметр с размерностью
энергии, который по обычаю также обозначают строчной \(h\).
Здесь мы сохраняем оба множителя, чтобы показать их физическое происхождение, а \(E\) — это
полная энергия.
\lecturetimestamp{00:48:21--00:49:27}}
\end{classroomqa}

\subsection{Подсчет конфигураций}

При фиксированных \(n\) и \(m=N-n\) нужно выбрать, какие \(n\) из \(N\) узлов направлены вверх.
Кратность равна
\begin{equation}
 \Omega(n,m)=\binom{N}{n}
 =\frac{N!}{n!\,m!}.
 \label{eq:sm8-spin-multiplicity}
\end{equation}
Каждая из этих конфигураций имеет одну и ту же энергию
\eqref{eq:sm8-counting-energy}.

\begin{classroomqa}
\audiencequestion{Могут ли \(n\) и \(m\) оба быть индексами суммирования, раз
они не независимы? И где появляется комбинаторный множитель?}
\lecturerresponse{Их можно записывать как связанную пару с
\(n+m=N\), либо можно суммировать по \(n=0,\ldots,N\) и положить \(m=N-n\).
Больцмановский вес относится к каждой отдельной конфигурации; при объединении всех
конфигураций с одной и той же энергией этот вес умножается на
\(N!/(n!\,m!)\).
\lecturetimestamp{00:52:04--00:53:03}}
\end{classroomqa}

Следовательно, статистическая сумма равна
\begin{align}
 Z_N
 &=\sum_{n+m=N}
   \frac{N!}{n!\,m!}
   e^{-\beta\mu H(n-m)}
 \notag\\
 &=\sum_{n=0}^{N}
   \binom{N}{n}
   \left(e^{-\beta\mu H}\right)^n
   \left(e^{+\beta\mu H}\right)^{N-n}.
 \label{eq:sm8-spin-z-sum}
\end{align}
Введем
\begin{equation}
 x=e^{-\beta\mu H},
 \qquad
 y=e^{+\beta\mu H}.
 \label{eq:sm8-xy}
\end{equation}
Затем вычисление завершается по биному Ньютона:
\begin{align}
 Z_N
 &=\sum_{n=0}^{N}\binom{N}{n}x^n y^{N-n}
   =(x+y)^N
 \notag\\
 &=\left(e^{-\beta\mu H}+e^{+\beta\mu H}\right)^N
   =\left[2\cosh(\beta\mu H)\right]^N.
 \label{eq:sm8-spin-z-closed}
\end{align}

\begin{figure}[t]
\centering
\includegraphics[width=0.90\textwidth]{lecture_08_figure_04.png}
\caption{Сумма для статистической суммы, подстановки \(x\) и \(y\) и
сведение по биному к степени \(2\cosh(\beta\mu H)\).}
\lectureframe{00:56:00}
\end{figure}

Есть и более короткий вывод. Поскольку спины независимы, для одного
спина имеем
\begin{equation}
 z_1=e^{-\beta\mu H}+e^{+\beta\mu H}
    =2\cosh(\beta\mu H),
 \label{eq:sm8-one-spin-z}
\end{equation}
и \(Z_N=z_1^N\). Тем не менее комбинаторный вывод стоит
сохранить, поскольку он явно показывает, как кратность конкурирует с энергией.

\section{Намагниченность конфигурации и ансамблевая намагниченность}

Для одной конфигурации определим ее безразмерный избыток спинов на узел как
\begin{equation}
 M_{\rm conf}=\frac{n-m}{N}.
 \label{eq:sm8-config-magnetization}
\end{equation}
Он изменяется от \(+1\), когда все спины направлены вверх, до \(-1\), когда все спины направлены вниз. Используемая
здесь термодинамическая намагниченность есть его среднее по Больцману,
\begin{equation}
 M=\left\langle M_{\rm conf}\right\rangle
  =\left\langle\frac{n-m}{N}\right\rangle.
 \label{eq:sm8-average-magnetization}
\end{equation}
Множителю \(\mu H\) не место в этом определении. Во время вычисления его ненадолго
туда вписали, а затем убрали; \(\mu H\) переводит избыток спинов в энергию.

Для каждой конфигурации по отдельности
\begin{equation}
 E=N\mu H M_{\rm conf}.
 \label{eq:sm8-energy-config-m}
\end{equation}
Линейность усреднения сразу дает
\begin{equation}
 \langle E\rangle=N\mu H M.
 \label{eq:sm8-energy-average-m}
\end{equation}

\begin{figure}[t]
\centering
\includegraphics[width=0.90\textwidth]{lecture_08_figure_05.png}
\caption{Ряд спинов, исправленное среднее значение намагниченности и соотношение с энергией,
подчеркнутое на доске.}
\lectureframe{01:01:22}
\end{figure}

\begin{classroomqa}
\audiencequestion{Приписывается ли намагниченность конкретной
конфигурации, или речь идет о средней намагниченности? Почему то же самое
соотношение справедливо и для средней энергии?}
\lecturerresponse{Здесь появляются оба понятия, и их следует различать.
\((n-m)/N\) — это значение в одной конфигурации; \(M\) — его среднее по Больцману.
Поскольку \(E=N\mu H M_{\rm conf}\) выполняется отдельно в каждой
конфигурации, усреднение обеих частей с теми же вероятностями дает
\(\langle E\rangle=N\mu H M\). Фиксированные величины \(N\), \(\mu\) и
\(H\) выносятся за знак среднего.
\lecturetimestamp{00:59:54--01:03:40}}
\end{classroomqa}

\begin{classroomqa}
\audiencequestion{Как модель может ориентировать каждый спин вдоль оси север--юг, не вводя предпочтения между направлениями вверх и вниз?}
\lecturerresponse{Отсутствие востока и запада заложено в двухсостояний
модель; это ограничение допустимых степеней свободы. В рассматриваемой
модели с внешним полем \emph{есть} различие между направлениями вверх и вниз: при нашем соглашении
о знаках поле понижает энергию спинов, направленных вниз. Следовательно, у этой
модели нет симметрии вверх--вниз. Вводимая ниже модель Изинга без поля
восстановит эту симметрию.
\lecturetimestamp{01:03:41--01:04:53}}
\end{classroomqa}

\subsection{Энергия из статистической суммы}

Средняя энергия следует из обычного тождества
\begin{equation}
 \langle E\rangle
 =-\frac{\partial}{\partial\beta}\log Z_N.
 \label{eq:sm8-energy-identity}
\end{equation}
Используя \eqref{eq:sm8-spin-z-closed},
\begin{align}
 \log Z_N
 &=N\log 2+N\log\!\left[\cosh(\beta\mu H)\right],
 \notag\\
 \langle E\rangle
 &=-N\mu H\,
   \frac{\sinh(\beta\mu H)}{\cosh(\beta\mu H)}
 \notag\\
 &=-N\mu H\tanh(\beta\mu H).
 \label{eq:sm8-average-energy}
\end{align}
Сопоставляя это с \eqref{eq:sm8-energy-average-m}, получаем
\begin{equation}
 \boxed{M=-\tanh(\beta\mu H)}.
 \label{eq:sm8-magnetization-tanh}
\end{equation}

Физический смысл заключен уже в виде ответа. Напомним,
\begin{equation}
 \cosh x=\frac{e^x+e^{-x}}{2},
 \qquad
 \sinh x=\frac{e^x-e^{-x}}{2}.
 \label{eq:sm8-hyperbolic-definitions}
\end{equation}
Вблизи нуля \(\tanh x=x+O(x^3)\); при больших положительных \(x\) он стремится к
\(1\). Поскольку \(\beta=1/T\),
\begin{equation}
 M\longrightarrow -1
 \quad (T\longrightarrow0),
 \qquad
 M\longrightarrow 0
 \quad (T\longrightarrow\infty).
 \label{eq:sm8-magnetization-limits}
\end{equation}
При нулевой температуре каждый спин находится в состоянии вниз с меньшей энергией. При
бесконечной температуре все конфигурации получают одинаковый вес, и предпочтение,
создаваемое полем, размывается. Изменение знака \(H\) меняет знак \(M\).

\begin{figure}[t]
\centering
\includegraphics[width=0.90\textwidth]{lecture_08_figure_09.png}
\caption{Кривая гиперболического тангенса стремится к постоянной при больших
\(\beta\mu H\) и линейна вблизи начала координат. Физическая намагниченность в
соглашении, принятом в лекции, равна отрицательной величине изображенной функции \(\tanh\).}
\lectureframe{01:13:14}
\end{figure}

\begin{classroomqa}
\audiencequestion{Намагниченность меняется непрерывно. Где же
температура фазового перехода, о которой упоминалось раньше?}
\lecturerresponse{В этой модели такой температуры нет. Независимые спины в
ненулевом внешнем поле слишком просты, и
\(\log Z_N=N\log[2\cosh(\beta\mu H)]\) гладка при любой конечной положительной
температуре. Для фазового перехода нужна взаимодействующая симметричная задача,
к которой мы теперь переходим; двумерная модель Изинга при нулевом поле — первый
представитель этого семейства, имеющий переход при конечной температуре.
\lecturetimestamp{01:14:47--01:15:11}}
\end{classroomqa}

\section{Модель Изинга: энергия в соседних парах}

Внешнее поле выбрало направление еще до начала статистического рассмотрения. Уберем
это поле и предоставим соседним спинам создавать друг для друга эффективное поле.
Если оба соседа данного узла направлены вверх, они благоприятствуют направлению вверх; если оба направлены вниз, они
благоприятствуют направлению вниз; если один направлен в одну сторону, а другой в другую,
то в простейшей картине их эффекты компенсируются. Теперь энергия связана с
парами, а не с отдельными узлами.

Для двух соседних спинов выберем \(J>0\) и запишем
\begin{equation}
 E_{12}=-J\sigma_1\sigma_2.
 \label{eq:sm8-pair-energy}
\end{equation}
Для сонаправленных спинов \(\sigma_1\sigma_2=+1\) и энергия равна \(-J\);
для противоположно направленных спинов произведение равно \(-1\), а энергия — \(+J\). Поэтому именно знак минус
делает взаимодействие ферромагнитным.

\begin{figure}[t]
\centering
\includegraphics[width=0.90\textwidth]{lecture_08_figure_06.png}
\caption{Локальное ферромагнитное взаимодействие
\(E_{12}=-J\sigma_1\sigma_2\).}
\lectureframe{01:19:59}
\end{figure}

Для открытой одномерной цепи гамильтониан имеет вид
\begin{equation}
 {\cal H}=-J\sum_{i=1}^{N-1}\sigma_i\sigma_{i+1}.
 \label{eq:sm8-ising-chain-open}
\end{equation}
При периодических граничных условиях полагают \(\sigma_{N+1}=\sigma_1\) и
суммируют до \(N\). В обоих соглашениях каждая пара соседних спинов учитывается один раз;
их различие на границе становится пренебрежимо малым для длинной цепи.

\begin{figure}[t]
\centering
\includegraphics[width=0.90\textwidth]{lecture_08_figure_10.png}
\caption{Локальная энергия пары, обобщённая до гамильтониана цепи с взаимодействием ближайших соседей.}
\lectureframe{01:20:46}
\end{figure}

При бесконечной температуре любая последовательность спинов равновероятна. В среднем в половине
связей спины сонаправлены, а в половине — противоположно направлены, поэтому
\begin{equation}
 \left\langle\sigma_i\sigma_{i+1}\right\rangle_{T=\infty}=0,
 \qquad
 \left\langle{\cal H}\right\rangle_{T=\infty}=0.
 \label{eq:sm8-infinite-t-bond}
\end{equation}
Нулевая энергия здесь не соответствует низкоэнергетическому состоянию. Полностью выровненное основное состояние имеет
энергию \(-J\) на связь, так что ансамбль при бесконечной температуре имеет относительно него
экстенсивный избыток энергии.

\begin{classroomqa}
\audiencequestion{Как среднее произведение может быть равно нулю, если каждое отдельное
произведение \(\sigma_i\sigma_{i+1}\) равно либо \(+1\), либо \(-1\)?}
\lecturerresponse{Среднее не обязано совпадать ни с одним из значений, принимаемых в отдельной
конфигурации. При бесконечной температуре сонаправленные и противоположно направленные пары соседних спинов встречаются
одинаково часто, поэтому их вклады \(+1\) и \(-1\) взаимно сокращаются. То же
различие между усреднением и отдельной конфигурацией, которое было важно для намагниченности, важно и здесь.
\lecturetimestamp{01:21:39--01:22:37}}
\end{classroomqa}

\subsection{Два основных состояния и одна глобальная симметрия}

При нулевой температуре в каждой связи спины стремятся быть сонаправленными. Поэтому имеются
два основных состояния,
\begin{equation}
 (+,+,+,\ldots,+)
 \qquad\hbox{and}\qquad
 (-,-,-,\ldots,-).
 \label{eq:sm8-two-ground-states}
\end{equation}
У них в точности одна и та же энергия. Гамильтониан делает это вырождение
наглядным: если инвертировать каждый спин,
\begin{equation}
 \sigma_i\longmapsto-\sigma_i
 \qquad\hbox{for all }i,
 \label{eq:sm8-global-flip}
\end{equation}
то произведение на каждой связи не изменится,
\begin{equation}
 (-\sigma_i)(-\sigma_{i+1})
 =\sigma_i\sigma_{i+1}.
 \label{eq:sm8-global-flip-invariance}
\end{equation}
Инверсия одного спина не является симметрией, потому что она меняет знак
произведений на соседних с ним связях. Инверсия всех спинов является симметрией. Это
глобальная симметрия \(\mathbb Z_2\) модели Изинга без поля.

\begin{figure}[t]
\centering
\includegraphics[width=0.90\textwidth]{lecture_08_figure_11.png}
\caption{Гамильтониан цепи рядом с глобальным преобразованием
\(\sigma_i\mapsto-\sigma_i\), которое оставляет инвариантным произведение каждой пары.}
\lectureframe{01:26:31}
\end{figure}

Приложим сколь угодно малое поле к одному спину при строго нулевой температуре. Оно снимает
вырождение между двумя основными состояниями, а ферромагнитные связи распространяют
этот выбор по цепи. Если убрать поле после того, как состояние уже
выбрано, большая система может очень долго сохранять свою ориентацию.
Это интуитивное введение в спонтанное нарушение симметрии.

Здесь есть одна существенная статистическая оговорка. Для конечной системы в
точно нулевом поле канонический ансамбль приписывает двум основным состояниям
одинаковые веса и, следовательно, нулевую среднюю намагниченность. Фаза с нарушенной симметрией
определяется тем, что термодинамический предел берут до снятия выбирающего
поля:
\begin{equation}
 M_*=\lim_{H\to0^+}\lim_{N\to\infty}M_N(H).
 \label{eq:sm8-order-of-limits}
\end{equation}
В одномерной модели с взаимодействием ближайших соседей этот предел отличен от нуля только при
\(T=0\); любая конечная температура разрушает дальний порядок в бесконечной
цепи. В двух измерениях упорядоченная фаза сохраняется ниже ненулевой
критической температуры.

\subsection{Историческая поправка}

В лекционном изложении эту историю рассказывают с шуткой за счёт Изинга.
Точная хронология доброжелательнее. Вильгельм Ленц предложил модель, а его
студент Эрнст Изинг решил одномерный случай в своей диссертации 1924 года,
опубликованной в 1925 году. Изинг правильно установил отсутствие ферромагнитного
перехода при конечной температуре в одном измерении, хотя более широкий смысл
более высоких размерностей тогда ещё не был понят. Крамерс и Ваннье разработали
двумерную дуальность и определили критическую точку квадратной решётки
в 1941 году; Онсагер дал знаменитое точное решение при нулевом поле в 1944 году.

Теперь ясно, каким должно быть следующее вычисление. Мы должны найти статистическую сумму для
\eqref{eq:sm8-ising-chain-open} и определить, как корреляции в этой модели зависят от
температуры. Это вычисление явно покажет, почему в одномерной
цепи нет фазового перехода при конечной температуре и почему размерность
имеет значение.

\section{Где мы находимся}

Задача о возвратах и задача о спинах — это две версии одного и того же
статистического урока. В газе дефицит энтропии делает макросостояние
экспоненциально редким:
\begin{equation}
 P(A)=e^{-\Delta S},
 \qquad
 \tau_{\rm rec}\sim\tau_{\rm mix}e^{\Delta S}.
 \label{eq:sm8-summary-recurrence}
\end{equation}
Обратимая механика допускает возврат, тогда как огромное отношение фазовых
объёмов объясняет, почему в обычном опыте наблюдается только релаксация. Оставшаяся
проблема стрелы времени связана со специальным космологическим начальным условием, а не
с нарушением микроскопической обратимости.

Для независимых спинов в поле,
\begin{equation}
 Z_N=\left[2\cosh(\beta\mu H)\right]^N,
 \qquad
 M=-\tanh(\beta\mu H).
 \label{eq:sm8-summary-spins}
\end{equation}
Эта система явно смещена внешним полем и плавно меняется с температурой. Взаимодействие Изинга
\begin{equation}
 {\cal H}=-J\sum_i\sigma_i\sigma_{i+1}
 \label{eq:sm8-summary-ising}
\end{equation}
устраняет внешнее смещение, восстанавливает глобальную симметрию вверх--вниз и ставит
настоящий вопрос об упорядочении: может ли симметричная многочастичная система выбрать одно из
двух упорядоченных состояний при ненулевой температуре?
